\chapter*{Preface}
Tome~6 opens the second part of the series, the graduate toolbox. It builds the stage on which modern analysis and
probability happen. It starts with what numbers are made of: sets, the natural numbers, and the real line, which
earlier tomes took as given. Then come limits without coordinates: metric spaces, topological spaces, compactness
and Baire's theorem, which says that in a complete space ``most'' objects can be strange. The second half is measure
theory: how to give every reasonable set a length, how to integrate against such a length, how lengths of
products and of images behave, which spaces of functions this produces, when a measure has a density, and finally
measures on metric spaces, where weak convergence and the dimension of fractals live.

Tome~6 settles debts. Tome~0 took for granted the axioms of $\N$ and $\R$, the existence of the integral of a
continuous function, the intermediate value theorem, extreme values, Heine's theorem, the mean value theorem and the
basic properties of $\exp$, $\cos$, $\sin$ and $\pi$; all of them are proved here, in Chapters~1, 2, 3, 4 and~7.
Tome~3 admitted the existence of Lebesgue measure, of infinite sequences of independent random variables and the
Fubini--Tonelli theorem; Chapters~6, 8 and~11 prove them. Each promise of an earlier tome is kept where it was
announced, and the text says so when it happens.

It assumes only Tome~0. References such as ``Tome~3, Theorem~5.4'' point to earlier tomes, for comparison or for
examples; nothing here depends on them. Exercises are graded \textbf{$\star$} warm-up, \textbf{$\star\star$} for a
good afternoon, \textbf{$\star\star\star$} gems, and every exercise has a hint and a complete solution at the end of
its chapter. Each chapter has a companion Python notebook. Abstract as this tome is, a computer is a good companion:
the notebooks enumerate the rationals, compute with exact real numbers, iterate contractions to draw fractals and
rank web pages, test the digits of $\pi$ for normality, integrate in a hundred dimensions by Monte Carlo, blur images
by convolution and measure the dimension of fractals.

The rules of the earlier tomes hold here too. Every statement is proved, or, if quoted only for culture, labelled
as a remark that says where it is proved. The only things taken without proof are the axioms of set theory, listed
in Chapter~1; each use of the axiom of choice is pointed out. Every solution starts with \emph{How to find it}
before the \emph{Write-up}. ``Pause and think'' boxes ask you to guess before you read on.

The chapters follow the plan of the series. Chapter~1 builds the numbers; Chapters~2--5 are topology, from metric
spaces to Baire's theorem; Chapters~6--8 construct measures and the Lebesgue integral and prove Fubini's theorem and
the change of variables formula; Chapters~9--11 study $L^p$ spaces, densities and differentiation, and measures on
metric spaces.
